\chapter*{Conclusion Générale}
\addcontentsline{toc}{chapter}{Conclusion Générale}

Ce Projet de Fin d'Études, réalisé au sein de NextStep IT, avait pour objectif de concevoir et de réaliser une plateforme PaaS serverless multi-tenant, bâtie sur Kubernetes, permettant à des clients de déployer et d'exploiter des applications conteneurisées sans en gérer directement la complexité d'orchestration sous-jacente, avec une mise à l'échelle automatique jusqu'à zéro instance, une intégration native avec un système de messagerie événementielle, et une facturation reflétant l'usage réel des ressources.

La démarche suivie, structurée selon la méthodologie agile Scrum et organisée en quatre releases successives, a permis de construire cette plateforme de manière incrémentale et validée à chaque étape. La Release 0 a établi le socle Cloud Native du système : cluster Kubernetes isolé par tenant, exécution serverless via Knative Serving, messagerie événementielle via Kafka et Strimzi. La Release 1 a transformé ce socle en cœur applicatif cohérent, organisé par domaine de gestion : identité et contrôle d'accès, cycle de vie complet d'une application, messagerie événementielle exposée au client, et observabilité applicative. La Release 2 a ouvert la plateforme à son exploitation commerciale et opérationnelle : facturation à l'usage et paiement en ligne, gestion des comptes clients, des quotas et de l'audit, supervision du cluster et page de statut public. La Release 3, enfin, a consolidé la solidité technique de l'ensemble, par le durcissement de la sécurité et l'industrialisation des chaînes de déploiement continu.

Sur le plan technique, ce projet a permis de mettre en pratique, dans un contexte réel et contraint, un ensemble de technologies cloud native encore peu matures et en évolution rapide (Knative Serving, Knative Eventing, Strimzi) en les combinant à un backend Spring Boot orchestrant le cluster de manière programmatique via le client Fabric8, et à des interfaces React exploitant ces capacités au travers d'une API REST sécurisée. Il a également permis d'expérimenter concrètement les difficultés propres à l'intégration de briques d'infrastructure hétérogènes : cohérence entre l'état applicatif persisté et l'état réel du cluster, fiabilisation d'une chaîne d'intégration continue reposant sur la construction d'images sans démon privilégié, ou encore arbitrages entre sécurité idéale et contraintes de délai et de modèle produit.

Ce mémoire s'est attaché, à chaque étape, à documenter avec la même rigueur les fonctionnalités abouties et les limites identifiées, qu'il s'agisse de choix assumés ou de chantiers inachevés faute de temps. Ces limites, loin de constituer un aveu d'échec, dessinent une feuille de route réaliste pour une éventuelle industrialisation ultérieure de la plateforme par NextStep IT : généralisation de l'infrastructure-as-code à l'ensemble du socle cluster, migration du flux d'authentification vers Authorization Code avec PKCE, isolation renforcée du cluster Kafka partagé entre tenants, intégration de tests automatisés et d'une analyse de vulnérabilités dans les pipelines, et refactorisation de la duplication de code entre les deux applications frontend.

Au terme de ce projet, la problématique posée en introduction — permettre à des clients multiples de déployer de manière autonome et sécurisée des applications conteneurisées, avec une mise à l'échelle automatique incluant le retour à zéro instance, une intégration événementielle native et une facturation à l'usage — trouve une réponse fonctionnelle complète et opérationnelle, dont les limites ont été explicitement circonscrites. Ce travail constitue ainsi, pour NextStep IT, une base concrète et documentée pour l'évolution de son offre Cloud vers des services de plateforme applicative à plus forte valeur ajoutée.
